\chapter{Familia y personas}
\section{Objetivos}
\begin{itemize}[leftmargin=*]
	\item Nombrar miembros de la familia y describir relaciones básicas.
	\item Usar posesivos y adjetivos simples para describir personas.
\end{itemize}

\section{Contenidos clave}
\begin{itemize}[leftmargin=*]
	\item Vocabulario: \textit{gezin/familie}, padres, hermanos, pareja, hijos.
	\item Posesivos: \textit{mijn, jouw, zijn, haar, ons/onze, jullie, hun}.
	\item Adjetivos comunes: alto/bajo, joven/viejo, simpático.
\end{itemize}

\section{Ruta del capítulo}
Seis secciones con vocabulario, posesivos y práctica guiada.

\section{Miembros de la familia}
\textbf{Vocabulario base.}
\begin{itemize}[leftmargin=*]
  \item Núcleo: \textit{ouders} (padres), \textit{vader/papa}, \textit{moeder/mama}, \textit{kind}, \textit{zoon}, \textit{dochter}.
  \item Hermanos: \textit{broer}, \textit{zus}; plural: \textit{broers}, \textit{zussen}.
  \item Abuelos y nietos: \textit{opa/oma}, \textit{grootouders}, \textit{kleinkind}.
\end{itemize}

\textbf{Mini práctica}
\begin{itemize}[leftmargin=*]
  \item Completa tu árbol familiar en 8 palabras neerlandesas.
  \item Di en voz alta quién vive contigo: \textit{Ik woon met \\_}. 
\end{itemize}

\section{Relaciones personales}
\textbf{Términos frecuentes.}
\begin{itemize}[leftmargin=*]
  \item Pareja: \textit{partner}, \textit{vriend/vriendin} (novio/a), \textit{man} (esposo), \textit{vrouw} (esposa).
  \item Otros: \textit{oom, tante, neef, nicht} (tío, tía, primo/sobrino, prima/sobrina según contexto).
  \item Hijastros y padrastros: \textit{stiefvader, stiefmoeder, stiefzoon, stiefdochter}.
\end{itemize}

\textbf{Mini práctica}
\begin{itemize}[leftmargin=*]
  \item Presenta a tu pareja o amigo: \textit{Dit is mijn \\underline{\hspace{1.2cm}}}.
  \item Escribe dos frases sobre tus tíos o primos.
\end{itemize}

\section{Edad}
\textbf{Preguntar y responder.}
\begin{itemize}[leftmargin=*]
  \item Pregunta: \textit{Hoe oud ben je?} / cortés: \textit{Hoe oud bent u?}
  \item Respuesta: \textit{Ik ben dertig jaar (oud).} Puedes omitir \textit{jaar} en informal.
\end{itemize}

\textbf{Números + años}
\begin{itemize}[leftmargin=*]
  \item Usa cardinales: \textit{achtentwintig}, \textit{vijfendertig}.
  \item Para terceros: \textit{Hij is tweeëntwintig}, \textit{Zij is veertig}.
\end{itemize}

\textbf{Mini práctica}
\begin{itemize}[leftmargin=*]
  \item Pregunta y responde edades con un compañero.
  \item Escribe tres frases sobre edades de tu familia.
\end{itemize}

\section{Estado civil}
\textbf{Vocabulario clave.}
\begin{itemize}[leftmargin=*]
  \item \textit{single} / \textit{alleenstaand} (soltero/a), \textit{getrouwd} (casado/a), \textit{gescheiden} (divorciado/a), \textit{weduwe} / \textit{weduwnaar} (viuda/o).
  \item En pareja sin matrimonio: \textit{samenwonend} (conviviente).
\end{itemize}

\textbf{Frases modelo}
\begin{itemize}[leftmargin=*]
  \item \textit{Ik ben getrouwd}. (Estoy casado/a.)
  \item \textit{Hij is gescheiden}. (Él está divorciado.)
  \item \textit{Wij wonen samen}. (Vivimos juntos.)
\end{itemize}

\textbf{Mini práctica}
\begin{itemize}[leftmargin=*]
  \item Completa: \textit{Mijn zus is \\_}. 
  \item Redacta dos frases describiendo estados civiles de conocidos.
\end{itemize}

\section{Descripción básica de personas}
\textbf{Adjetivos útiles.}
\begin{itemize}[leftmargin=*]
  \item Apariencia: \textit{lang, klein, groot, slank, stevig}.
  \item Edad y carácter: \textit{jong, oud, aardig, rustig, druk, grappig}.
  \item Cabello: \textit{lang, kort, blond, donker, krullend}.
\end{itemize}

\textbf{Estructuras}
\begin{itemize}[leftmargin=*]
  \item Predicativo: \textit{Mijn broer is lang}.
  \item Atributivo: \textit{Ik heb een kleine zus}.\ \textit{-e} en adjetivos delante de sustantivos con \textit{de} y la mayoría de \textit{het}.
\end{itemize}

\textbf{Mini práctica}
\begin{itemize}[leftmargin=*]
  \item Describe a dos familiares en 2 frases cada uno.
  \item Escoge tres adjetivos para ti y léelos en voz alta.
\end{itemize}

\section{Hablar de otras personas}
\textbf{Pronombres y posesivos.} Revisa \textit{zijn} (su de él) y \textit{haar} (su de ella); plurales: \textit{hun} (su de ellos).

\textbf{Frases modelo}
\begin{itemize}[leftmargin=*]
  \item \textit{Zijn moeder woont in Utrecht}. (Su madre vive en Utrecht.)
  \item \textit{Haar broer is dertig}. (Su hermano tiene treinta años.)
  \item \textit{Hun kinderen studeren}. (Sus hijos estudian.)
\end{itemize}

\textbf{Mini práctica}
\begin{itemize}[leftmargin=*]
  \item Escribe tres frases sobre familiares de un amigo usando posesivos correctos.
  \item Corrige una frase con posesivo incorrecto (crea tu propio ejemplo).
\end{itemize}


\section{Pronunciación}
\begin{itemize}[leftmargin=*]
	\item Diptongo en \textit{huis/gezin}; evita \textit{huis} como /jus/.
	\item Sonoriza la \textit{g} suave en \textit{gezin}; no \textit{gesin} con /s/.
\end{itemize}

\section{Errores comunes}
\begin{itemize}[leftmargin=*]
	\item Confundir \textit{gezin} (núcleo hogar) con \textit{familie} (familia amplia).
	\item Olvidar concordancia de \textit{ons/onze} según género de palabra.
	\item Colocar adjetivo sin \textit{-e} tras \textit{de}-woorden: \textit{de grote broer}.
\end{itemize}

\section{Práctica guiada}
\begin{itemize}[leftmargin=*]
	\item Completa un árbol familiar con posesivos apropiados.
	\item Ordena oraciones con adjetivos y posesivos.
	\item Rellena diálogos preguntando por estado civil y edades.
\end{itemize}

\section{Ejercicios de producción}
\begin{itemize}[leftmargin=*]
	\item Escribe y graba una presentación de tu familia (10 frases).
	\item Describe a dos personas (apariencia + relación) en 6 frases cada una.
\end{itemize}

\section{Mini repaso}
\begin{itemize}[leftmargin=*]
	\item Diferencias \textit{gezin} y \textit{familie} en contexto.
	\item Usas posesivos correctos y concordados.
	\item Describes miembros con 2–3 adjetivos básicos.
\end{itemize}

\section{Resultado esperado}
\begin{itemize}[leftmargin=*]
	\item Al terminar puedes presentar a tu familia, describirla con adjetivos simples y preguntar por la de otros con posesivos correctos.
\end{itemize}

\section{Práctica sugerida}
\begin{itemize}[leftmargin=*]
	\item Árbol genealógico sencillo y frases descriptivas.
	\item Descripción oral de fotos familiares.
	\item Preguntas/respuestas sobre familiares de un personaje ficticio.
\end{itemize}

\section{Microevaluación}
\begin{itemize}[leftmargin=*]
	\item Presenta a tu familia en 8 frases grabadas o escritas.
\end{itemize}
